\documentclass[11pt,a4paper]{article}
\input{preamble}
\begin{document}
\coverpage{Project Diary}{A Record of Decisions, Implementations and Corrections}

\newpage
\section*{About this Diary}
\addcontentsline{toc}{section}{About this Diary}
\markboth{About this Diary}{}

This diary records the development of ThinkStack from 5 June to 7 August 2026.
It is reconstructed from the repository itself: 250 commits, 59 merged pull
requests, 51 recorded architecture decisions, and the published releases. It is
not written from memory.

Two conventions are observed throughout. First, every quantity was measured or
counted, never estimated. Second, decisions that turned out to be wrong are
recorded alongside those that were right, because a diary that contains only
successes teaches nothing.

The work was carried out jointly by the three team members. Entries describe
what was decided and built rather than who performed each action.

\vspace{0.6cm}
\noindent\textbf{Team:} Aditya Mehta (2440204), Jitvan Chadha (2440222), Ritesh K R (2440233)

\newpage
\tableofcontents
\newpage

\section{Timeline at a Glance}
\markboth{Timeline}{}

\begin{center}
\small
\begin{tabularx}{\linewidth}{@{}llX@{}}
\toprule
\textbf{Date} & \textbf{Phase} & \textbf{Event} \\
\midrule
5 Jun & Sprint I & repository created \\
17 Jun & Sprint I & project renamed ScholarLens to ThinkStack; decision log started \\
19 Jun & Sprint I & desktop shell introduced \\
20 Jun & Sprint I & encryption implemented \\
23 Jun & Sprint I & LaTeX editor and auto-healing compiler \\
24 Jun & Sprint I & interface unified; Sprint I closes \\
7 Jul & Sprint II & download page \\
10 Jul & Sprint II & cross-platform builds and CI \\
21 Jul & Sprint II & hardware-aware model routing \\
22 Jul & Sprint II & in-application updates \\
25 Jul & Sprint II & test suite introduced; Rust startup diagnosis \\
27 Jul & Sprint II & first public release, v0.1.0 \\
29 Jul & Sprint II & v0.1.1, then v1.0.0 \\
30 Jul & Sprint II & v1.0.0 found not to start; root cause analysis and correction \\
30 Jul & Sprint II & TeX engine bundled; validated beta published \\
31 Jul & Sprint III & release pipeline consolidated; every dependency pinned \\
1 to 3 Aug & Sprint III & LitGraph: search, analysis and gaps become one canvas \\
1 Aug & Sprint III & stale interface after update traced to the webview cache \\
2 Aug & Sprint III & capability model extracted; Bench introduced \\
4 Aug & Sprint III & auto-update failure diagnosed; installs on a dead channel rescued \\
4 Aug & Sprint III & beta v1.9.0 published \\
4 to 5 Aug & Sprint III & model registry, task router, and Hugging Face acquisition \\
5 Aug & Sprint III & baseline chosen by measurement; render tests added \\
\bottomrule
\end{tabularx}
\end{center}

\section{Sprint I: Building the Application}

\subsection{5 to 16 June: foundations}

The repository was created on 5 June. The first two weeks established the
project structure, an initial working demonstration, and the first iteration of
the interface.

\subsection{17 June: naming, audit, and the decision log}

Three things happened on the same day that shaped everything afterwards.

The project was renamed from ScholarLens to ThinkStack. A backend audit
corrected a set of defects found by reading rather than by testing. Most
consequentially, an architecture decision log was started.

\evidence{\textbf{Decision recorded.} Run entirely on the user's device, with no
hosted language model. The reasoning: a researcher working with unpublished
results or institutionally restricted material cannot upload it, and a student
project has no budget for per-call inference. This single decision determined
local inference, quantised models, a frozen runtime, and eventually a bundled
typesetting engine.}

The log reached 32 entries by the end of the project. Its value was not
apparent at the time; it became apparent during the dependency audit in the
final week, when the reasoning behind earlier choices could be read rather than
reconstructed.

\subsection{19 to 21 June: desktop shell and encryption}

The application was wrapped in a desktop shell. The shell starts the backend as
a supervised child process and stops it when the window closes, so no orphaned
process is left holding a port.

Encryption was implemented with Argon2id for key derivation and AES-256-GCM for
the ciphertext.

\evidence{\textbf{Decision recorded.} Use authenticated encryption rather than
an unauthenticated mode. Tampering with a stored paper is then detected at
decryption instead of producing plausible but corrupted output, which for
research material is the difference between an error and a silent falsification.}

\subsection{23 to 24 June: the paper writer}

The LaTeX editor was implemented, with a compiler designed to be forgiving. A
small language model produces markup that is frequently almost correct, so the
compiler inserts missing package declarations, wraps bare fragments into
complete documents, and continues past recoverable errors. When a pass produces
no PDF, the failing environment is progressively neutralised and compilation is
retried, so one malformed figure does not cost the author the paper.

Sprint I closed on 24 June with every feature working from a source checkout.

\section{Sprint II: Making It Usable by Others}

\subsection{7 to 10 July: distribution}

A download page was produced, and cross-platform builds were automated. The
build configuration was made data-driven, so that adding an operating system is
an edit to a configuration file rather than to build logic.

\subsection{20 to 22 July: hardware awareness}

Model selection was made dependent on the machine rather than fixed.

\evidence{\textbf{Decision recorded.} Size the model against \emph{available}
memory rather than installed memory. A user has a browser and an editor open.
Sizing against total memory produces a process that is terminated by the
operating system partway through an analysis, which is worse than producing a
slightly weaker answer.}

A related decision concerned graphics hardware. The diagnosis reports the GPU
but enables offload only for one specific class of device.

\evidence{\textbf{Decision recorded.} Report the GPU, but treat all but CUDA
devices with sufficient memory as unusable. The inference library shipped in the
installer is a CPU build, so attempting offload elsewhere fails at model load.
Reporting a capability the runtime cannot use converts a working configuration
into a crash.}

\subsection{25 July: testing and the Rust rewrite}

An automated test suite was introduced. Startup diagnosis was moved from Python
into the desktop shell, written in Rust, because the Python implementation
imported a deep learning framework purely to determine whether a GPU existed,
which cost several seconds on every launch.

\evidence{\textbf{Decision recorded, with a limit.} Hardware diagnosis moved to
Rust because it was on the startup path and slow. Model discovery was measured
at approximately 9 milliseconds and was left in Python, because the same
argument does not apply to code that is already fast. The principle recorded was
to move work for a measured reason, not because a language seems faster.}

\subsection{27 to 29 July: first releases}

Version 0.1.0 was published on 27 July, the first release to build on all three
operating systems. Two defects had blocked it, both worth recording because
neither was reproducible locally.

\begin{itemize}[leftmargin=*]
  \item The Windows build failed at model download. A JSON tool on that platform
        emits carriage returns, which left an invalid character on the end of a
        URL. The transfer tool rejected it. Linux and macOS were unaffected, so
        the defect existed only on the platform none of the team used daily.
  \item The macOS build failed at code signing. An absent optional credential
        was passed as an empty string, so the signing tool attempted to import
        an empty certificate rather than skipping signing.
\end{itemize}

Version 1.0.0 followed on 29 July, alongside a decision to reduce the installer
size.

\evidence{\textbf{Decision recorded.} Ship only the smaller model and offer the
larger one with the user's consent, reusing weights the user may already have
through other local tools. This removed approximately one gigabyte from every
installer. Matching required a canonical naming key, because the same weights
are named differently by each runtime, and comparing filenames caused the
application to offer a download the user did not need.}

\subsection{30 July: the release that did not work}

\evidence{\textbf{Observed.} The published v1.0.0 installer, downloaded and run,
displayed a loading screen indefinitely. One tester waited approximately 200
seconds. Continuous integration had been green throughout.}

This is the most instructive day of the project.

\subsubsection{Diagnosis}

The startup path was instrumented before anything was changed, on the reasoning
that the failure could not be fixed while it was invisible. The instrumentation
identified the cause in a single run.

Three defects compounded:

\begin{enumerate}[leftmargin=*]
  \item The packaging framework installs the executable and its resources into
        different directories. The lookup searched only beside the executable,
        found nothing, and fell back to the development path of launching a
        system Python interpreter.
  \item The Linux application image exports an environment variable that
        redirects a Python interpreter's library search into the mounted image.
        A frozen interpreter ignores this; a system interpreter does not, and it
        terminated before running any project code.
  \item The loading screen retried without limit and had no error state, so a
        backend that died in under a second looked exactly like one that was
        slow.
\end{enumerate}

\subsubsection{What the same investigation uncovered}

\begin{itemize}[leftmargin=*]
  \item The embedding weights had never been placed in the bundle. The first
        document ingested on an installed copy attempted a network download, in
        an application whose entire premise is that it needs no network. The
        code's own documentation stated the opposite.
  \item The model directory was derived from the current working directory,
        which the application image changes before launching. This was incorrect
        on every operating system for every installed copy; it worked only when
        launched from a developer checkout.
  \item The window terminated approximately one second after appearing, from
        heap corruption in a graphics path of the system webview. Every backend
        check passed while this happened, because the backend was healthy and
        only the window was gone.
\end{itemize}

\subsubsection{The correction that mattered most}

Rather than continue fixing defects individually, every external dependency was
enumerated by walking all imports and every subprocess call across 51 backend
modules.

\evidence{\textbf{Result.} Exactly one unbundled runtime dependency existed: the
TeX engine. Everything else was already inside the bundle. Bundling that engine
closed the entire category rather than one instance of it.}

The pattern the audit exposed was that the absent TeX engine, the absent
embedding weights and the incorrect model path were the same defect: an
assumption that a user's machine resembles a developer's.

\subsection{30 July: bundling the typesetting engine}

The engine was measured before it was adopted: a 58 MB binary and a package
cache of approximately 47 MB, against roughly 1.1 GB of headroom under the
hosting platform's asset limit.

The first attempt worked locally and failed on the Linux build server, which
produced the most transferable lesson of the project.

\evidence{\textbf{Cause.} The chosen engine version required a newer system C
library than the build platform provides, so it never executed at all, failing
in 15 milliseconds. The build step recorded a warning and continued, so the
installer was assembled with an engine that had no package cache. The same
binary would have reached users of long-term-support distributions, who would
have had a paper writer that silently could not compile.}

\evidence{\textbf{Correction.} The engine version was chosen by measuring the
minimum system library each release requires, rather than by taking the newest
release. The selected version runs on the build platform and on older user
systems, and reduced the bundled payload from 104 MB to 70 MB. Equally
important, the build step now prints the engine's actual error and stops the
build, because suppressing that error is what made the diagnosis expensive.}

\subsection{30 July: verified delivery}

A validated beta was published to all three operating systems, with every
automated check passing, including compilation of a PDF by the packaged
application.

\section{Sprint III: Correctness of What Was Delivered}

Sprint II ended with software that installed. Sprint III began with the
discovery that installing is not the same as remaining correct, and most of its
work was caused by faults that only appear after a second release exists.

\subsection{31 July to 1 August: the pipeline, and what a green tick means}

Five release workflows had accumulated, differing only in what started them and
how the version was derived; their build and publish halves were identical.
They were consolidated into one, with the channel decided by which branch
received the merge.

Two failures during this period shared a cause. A macOS launch check had been
ordered before the macOS build and marked \texttt{continue-on-error}, so it
reported success for a run in which it never executed. Separately, an
unpinned transitive dependency updated overnight and broke all three platforms
from a commit that changed no application code. Both were cases of a signal
that appeared to mean verification and did not.

A third fault was subtler. Users who updated correctly continued to see the
previous interface, because the embedded webview served \texttt{index.html}
from its own cache, which outlives the application. The correction was to serve
the entry document with \texttt{no-store} while allowing content-hashed assets
to remain immutable.

\subsection{1 to 3 August: LitGraph}

The gap finder, search, and summarisation had been separate screens over the
same corpus and much of the same computation. They were consolidated into one
canvas, in which the literature is presented as a connected structure rather
than as three lists: a question is asked once, the passages that answer it are
shown, the papers containing them are shown in relation to one another, and the
places where the literature is thin are marked on the same surface.

Analysis was moved off the request path and precomputed at ingestion, because a
canvas that must summarise before it can draw is a canvas that appears broken.
Selecting a node opens the paper itself, with the passages the map relied on
marked in it, so that a claim on the canvas can be checked against its source
without leaving the screen.

This was the largest single feature of the project and closes the consolidation
that had been recorded as planned work since Sprint I.

\subsection{2 August: separating measurement from decision}

Hardware classification existed twice: the Rust shell measured the machine at
startup, and the Python backend measured it again. The two disagreed, and the
Rust answer silently won, leaving the Python path authoritative-looking and
dead.

The responsibilities were separated. Rust measures; a single Python module
converts those measurements into decisions: context size, offload, token
budgets, and whether a document must be summarised in several passes. One
consequence was immediate: graphics offload had been decided by a test for CUDA
and dedicated video memory, which Apple Silicon can never satisfy despite
having a usable GPU, so every Mac had been running on the processor alone.

The same reorganisation gave the interface a place to show its findings, and
\emph{Bench} was introduced beside Library, LitGraph and Scribe: what the
machine can run, and what it runs it with.

\subsection{4 August: an update mechanism that could not succeed}

A tester reported that the update function did nothing on macOS. The
investigation found three faults in sequence, each concealing the next.

The update endpoint is compiled into the application and cannot be changed
afterwards. A rolling release tag had been renamed some days earlier, so every
installation built during that window continued to poll an address that still
existed and still advertised the build they already had. Separately, the
manifest carried a pre-release version string while the application reported
the plain version, and semantic versioning ranks a pre-release below its own
release, so even a correctly aimed installation was told it was current.

The stranded installations could not be reached by any change to the source,
because the address they poll is fixed in the binary they are running. They
were recovered by publishing a manifest at the abandoned address that advertises
the current version and points at the current files, which is a redirection in
everything but name. Testers on the beta channel now update in place.

The lesson generalised: the earlier failure and this one both arose from a
change that did not declare what it superseded, leaving the system to infer
intent from absence.

\subsection{4 to 5 August: models the user chooses}

Until this point the application used the model it shipped with, and one
optional larger model, selected by a fixed table of filenames. A researcher who
already had suitable weights, through Ollama, LM Studio, or a previous
download, could not tell the application to use them.

Three modules were added. A \emph{registry} records the models available to an
installation and which task each may perform, persisted through the existing
crash-safe write. A \emph{router} answers which model serves a given task on
this machine at this moment, with every dependency supplied by its caller, so
that routing can be tested against fabricated hardware without a model runtime
present. A \emph{manifest}, generated during the build from the release
configuration, records what that build bundled and which previously bundled
models it replaces, so that exchanging the included model becomes one
configuration edit rather than four coordinated source changes.

Imported models are referenced where they already are, never copied: the
alternative would double the disk cost of a file measured in gigabytes on
machines already selected for being constrained. A separate flag records
whether the application created a file, and nothing deletes a file it did not
create.

Acquisition from Hugging Face was added last, as the only part of the
application that reaches the network. It searches on submission rather than on
keystroke, states plainly that it is doing so, and constructs its download
address from a repository identifier and a filename rather than accepting an
address, since the local interface is reachable from the embedded browser.

Three defects found during this work shared a single form. A size was rounded
before being compared against a memory budget, so any file below roughly five
megabytes measured as zero and satisfied every constraint. A measured size of
zero was treated by a default-value expression as \emph{no measurement}, so a
stale declared figure of ninety-nine gigabytes was preferred over the truth.
And the memory budget itself, computed as free memory less a fixed reserve,
floors at zero on a constrained machine: where zero already meant
\emph{unmeasured, therefore unconstrained}, so the machine least able to run
anything was told it could run everything.

In each case a legitimate zero was indistinguishable from an absent value. The
correction in each was the same: ask the module that owns the decision, which
compares against real quantities, rather than testing a number for truthiness.


\section{Decisions and Their Consequences}

\begin{center}
\small
\begin{tabularx}{\linewidth}{@{}XXX@{}}
\toprule
\textbf{Decision} & \textbf{Reason} & \textbf{Consequence accepted} \\
\midrule
local inference only & privacy, cost, offline operation & far weaker models; seconds per response \\
quantised small models & fit consumer memory & reduced quality on structured output \\
frozen Python runtime & the user installs nothing & installers approach 1 GB \\
file-based vector store & no database to install & keyword index rebuilt per query \\
bundle only the smaller model & installer size & analysis quality depends on a consented download \\
bundle a typesetting engine & the writer must work with no LaTeX & 70 MB added per installer \\
CPU-only inference build & one build behaves identically everywhere & GPU hardware unused \\
tags, not branches, produce releases & releases cannot happen by accident & a deliberate step is required \\
updates only when requested & an offline product must not contact a server unprompted & a user who never asks never updates \\
merges produce releases, tags record them & a stray tag must not publish to every user & releasing requires a reviewed merge \\
imported models referenced, never copied & a multi-gigabyte file must not cost twice the disk & the application depends on a path it does not own \\
task assignment chosen by the user & parameter count does not predict structured output & one extra step when adding a model \\
download addresses constructed, never accepted & the local interface is reachable from the embedded browser & only known hosts can be fetched from \\
\bottomrule
\end{tabularx}
\end{center}

\section{What Was Learned}

\subsection{Testing source code is not testing a product}

309 tests passed while the shipped installer could not start. The tests examined
a source tree; users run an installer. The verification tiers added afterwards
exercise the packaged application on each supported operating system, on
machines that have no model and no typesetting distribution installed, so that a
pass is evidence about a user's machine.

\subsection{Suppressed errors are expensive}

Two separate diagnoses were slowed by output sent to a null device. In both
cases the fix was to print the error and stop, rather than warn and continue.
A build that continues after a step has failed produces an artefact nobody
intended to ship.

\subsection{Fix categories, not instances}

Three defects were the same defect. Enumerating every dependency once bounded
the problem and ended it, where fixing them one at a time as testers reported
them would not have converged.

\subsection{Documentation drifts silently}

Instructions told macOS users to use a mechanism the operating system had
removed, and a LaTeX error message offered a command for one Linux distribution
to every user on every platform. Both were written when they were true. Neither
was noticed until a person followed them and could not proceed.


\subsection{5 August: choosing a model by measuring it}

The bundled model was to be replaced with something lighter and newer. Qwen3
0.6B was the obvious candidate: smaller than the model it would replace, a
generation ahead, and its download verified. Every structural check passed.

It returned an empty object.

Gap finding constrains the model's output with a grammar that permits only JSON
from the first token. A reasoning model begins by thinking, and the grammar
leaves that nowhere to go, so it emits the shortest legal document instead:
\texttt{\{\}}. The parser succeeds. Nothing downstream notices. The feature
returns no gaps and reports no error.

A bake-off was written against the two flagship jobs, and it changed the
decision entirely. The incumbent Qwen2.5 0.5B recovered four of six stated
limitations, produced valid LaTeX, and was the fastest of the four candidates
tested; nothing larger justified its additional weight. The model that had been
about to ship on the strength of its specification came last.

One further correction belongs here, because it was the harness rather than the
product. The LaTeX test initially wrapped its output in JSON to make it
machine-checkable, and two candidates were marked down for producing corrupt
markup: inside a JSON string, \verb|\begin| is a valid escape and the parser
silently converts it to a control character. Scribe does not do this; it asks
for plain text. A test that differs from the production path measures the test.

\subsection{5 August: a test that had never failed}

Three times in one week the application shipped a blank page: the build
succeeds, because the markup is valid, and the interface framework then throws
at render time on a name that no longer exists, leaving an empty document and
an error visible only in a browser console.

A test was added to mount every screen. It passed. The bug was then reproduced
deliberately, the import removed again, and the test still passed, because
it mounted the four page components and never the application shell, which is
where the missing import was. It would have passed indefinitely while the
product failed to start.

The same exercise was applied to the rest of the new work: each guarantee was
broken on purpose to confirm something noticed. Two of five did not. The
protection against deleting a file the application did not create had a clear
error message, an obvious intent, and no test at all.

The general lesson is worth stating plainly, because this project has now
learned it twice in different forms. A suite that has never failed is a
hypothesis. Sprint II established that testing source code is not testing a
product; this establishes that writing a test is not the same as being
protected by one.


\subsection{6 August: a default that was indistinguishable from a decision}

Before beginning work on the paper editor, the repository was surveyed for
material that no longer had a purpose. Two modules in the gap-finding package
were found to be unreachable: an earlier implementation that asked the model for
gaps and for research directions in two separate calls, replaced some weeks
earlier by a single merged call. A configuration file for a hosting provider the
project no longer deploys to was removed likewise, and a file named as though it
belonged to the automated suite was renamed, since the test runner is scoped to
one directory and had therefore never collected it. Roughly two hundred lines
were deleted and nothing referred to them.

Reading the surviving implementation next to the one it replaced revealed
something the deletion had nearly concealed. The older module named the task it
was performing when it called the language model; the merged replacement did
not. That argument is what selects the model. The application allows a user to
assign a particular model to gap finding, and the routing layer honours that
assignment by looking up the task it is given. Given no task, the argument takes
its default value of \emph{general}, so the lookup is performed against the
wrong entry, matches nothing, and falls back to the bundled model: while the
interface continues to display the model the user selected. Gap finding, one of
the two features the project presents as its principal contribution, had been
running on the smaller model and disregarding the user's choice entirely.

The fault is the same one recorded on 4 August for the writing task, and it
arrived by the same route: two calls were merged into one and the routing
argument was not carried across. What makes this class of defect durable is that
its symptom is indistinguishable from correct operation. Nothing fails. The
output is plausible, because a smaller model still produces gaps; it produces
worse ones. The only external evidence is a quality difference that a reader
would attribute to the difficulty of the task.

Following the principle recorded earlier in this diary: that categories
should be fixed rather than instances, the correction was accompanied by two
forms of protection. The first drives each feature with the model replaced by a
recording stand-in and asserts on the task the feature \emph{requested}, rather
than on the answer it received. The second parses the source of every module and
fails if any call to the generation interface omits the argument at all. The
second was written because the first can only cover entry points somebody
remembered to enumerate.

Run once, it reported two further instances that had not been suspected:
per-document analysis, merged from two passes in exactly the same way and losing
the argument in exactly the same way, and metadata extraction at ingestion. The
second proved not to be a fault. Extracting a title and an author list is
shallow work performed on every upload, with a regular-expression fallback
immediately beneath it, and directing it to the larger model would make every
ingestion wait on a model exchange for no benefit. It was therefore made to
state \emph{general} explicitly. This is the more useful half of the outcome:
the remedy is not that every call must name a heavier model, but that every call
must record which model it intends, so that a deliberate choice and a forgotten
argument stop appearing identical in the source.

The correction was verified by restoring the defect deliberately and confirming
that both protections failed, then removing it again and confirming that both
passed.

A feature that had been abandoned in planning was also removed in full. A
conversational assistant remained present as a mounted interface with no way to
reach it, and was still described as working in the user-facing documentation.
Removing it exposed a dependency worth recording: the paper editor imported its
retrieval step from that module, because assembling relevant excerpts within a
size budget had been written there first. That step is not conversational: it
takes a query and returns text, so rather than delete it, it was moved beside
the search implementation it is built upon and given a name describing what it
does. Deleting the feature without reading its dependents would have removed
the grounding from the paper editor, which is the feature the following work was
about to address.


\subsection{6 August: a passing suite for a feature that did not work}

The sidebar was made to withdraw when a user begins working, giving the width
to whichever screen they are on, and to return when the logo is pressed. The
requirement as first implemented was read as a list: opening a paper on the map,
opening a document in the editor, expanding an entry in the library, and any
operation on the model bench. Each of those handlers was made to ask the shell
for room, and each was covered by a test.

The tests passed. The feature was reported as not working on two of the four
screens. Both accounts were correct. Everything on the list behaved; everything
absent from it did nothing, and most of what a person actually presses on those
two screens was absent. The specification had been satisfied and the requirement
had not, because the requirement was a property, any interaction, and it
had been implemented as an enumeration.

The correction removed more code than it added. The shell now observes its own
content region and withdraws on any press within it, so no feature calls
anything, imports anything, or can forget to. Controls nobody listed are covered,
including those not yet written. Two details were necessary rather than
decorative: the observation is of the press rather than the click, because
selecting a region of the map is a drag that may never produce a click; and it
is taken during the capture phase, because several rows deliberately stop
propagation on their own buttons and a listener waiting for the event to bubble
would never be told.

A second distinction was needed once actions could close the sidebar. Its state
had been remembered between runs, correctly, since a person who collapses it has
expressed a preference. An automatic collapse writing to that same record would
mean opening a single paper teaches the application to start collapsed
permanently, a setting the user never chose, altered by an action unrelated to
settings, with nothing on screen to explain it. The deliberate and the automatic
are now separate, only the first is remembered, and the logo clears both, since
clearing only its own would leave it apparently dead while a feature still held
the other.

The lesson is a narrower form of one this diary has already recorded twice. It
has been noted that a suite which has never failed is a hypothesis, and that
categories should be corrected rather than instances. This adds that a test can
fail in neither of those ways and still be worthless: it can verify precisely
what was implemented while the implementation answers a different question from
the one asked. Tests written from a list inherit the list's omissions, and they
inherit them silently, since a case nobody thought of is a case nobody writes an
assertion for.

\subsection{7 August: a project was always a directory}

A user reported that a colleague's paper would not compile. The engine said it
could not load a figure, and then that it had divided by zero. The second
message was the first one's consequence: the graphics package had read a width
from a file it never opened.

The cause was not in the compiler. A paper had been stored as a directory since
the first version, and compilation had always run with that directory as its
working directory, so the relative path in the document was correct. What was
missing was any way to put a second file into the folder. The document was the
only file the application could reach, so a figure could be referenced and never
supplied.

The correction was to expose the directory that already existed: list, read,
write, upload, rename, move, copy and delete, presented as a file tree with the
papers themselves as its top level. The row of chips that had previously
selected a paper was removed: it was a second picker for the upper level of
the same hierarchy, and there had been no way at all to reach the level beneath
it.

The security work is the part worth recording. The interface reaching this API
is a web view, which the desktop framework treats as remote content, so a
filename arriving here is untrusted input in the same sense that a downloaded
address is: \texttt{../../../.ssh/id\_rsa} is a filename. Every operation
therefore resolves its argument against the project directory and refuses
anything that lands outside. Resolution happens \emph{before} the check rather
than after, which is what catches a symbolic link: nothing about the name
\texttt{notes.tex} suggests it points at the system password file, so an
inspection of the string would pass it through. Deleting the containment check
was tried deliberately, and the tests fail.

A second failure arrived the same day from the same paper. Documents with an
index compiled to \texttt{Undefined control sequence \textbackslash indexentry}.
An index requires two passes with an external step between them: the first pass
writes a list of entries, a program converts that list into a formatted index,
and the second pass includes it. The bundled engine performs the equivalent step
for bibliographies automatically but not for indexes, so the intermediate file
was never converted and the index package fell back to reading the raw list as
if it were markup.

The obvious remedy was to call the standard index program. It is present on the
developer's machine and would not be present in the installer, which is the same
error as assuming a system typesetting installation: precisely what bundling
an engine was meant to eliminate. The index is therefore generated by the
application itself, which is a modest amount of code for a small and stable
format, and is consistent with the compiler's existing habit of repairing what
it is given rather than refusing it.

\subsection{7 August: interface controls that a desktop application does not have}

Three menu items, new file, new folder, new paper, were reported as doing
nothing at all. They had been built on the web browser's prompt dialog, which
the desktop web view does not implement; it returns nothing, and the code read
that as the user having cancelled. Deletion had the same dependency on the
browser's confirmation dialog, where the failure would have been worse: a guard
written as ``if not confirmed, stop'' either prevents the action permanently or,
read the other way, performs it without ever asking.

Naming now happens in a row within the tree itself, where the file will appear,
and destructive actions use a dialog belonging to the application. That dialog
was made shared rather than local, because it was about to be the third
independent implementation in the codebase; three dialogs is three sets of
keyboard handling, focus behaviour and dismissal rules to get right, and the
safest button is the one that does nothing, which is where focus now begins.

Behind that fault was a second one that would have prevented the menu from
working even with a functioning dialog. The menu dismissed itself on any press
elsewhere, and listened for that press during the capture phase: necessary,
because much of the tree stops events propagating. It did not check whether the
press was inside the menu, so pressing an item removed the button before the
click could reach it. Neither defect was visible from the code; both required
the application to be run.

The recurring lesson is the one this project has now recorded in several forms.
The suite was green, the components rendered, and the feature did not exist.
What was missing was not a test but an execution: none of this can be discovered
without launching the thing a user launches.

\subsection{7 August: the constraint was packaging, not hardware}

A tester reported that the application told him his graphics card could not be
used. The message was accurate. The inference engine is shipped compiled for the
processor only, so the card genuinely could not be used, and the diagnosis had
correctly said so. What the message did not do was offer anything, explain why,
or indicate whether it would ever change. A correct sentence at a dead end.

The first design took the obvious route. The manufacturer of that card publishes
prebuilt libraries, so the plan was to download those on request. It was
implemented to the point of working before a more careful look at the numbers
undermined it. Those libraries are approximately four hundred megabytes, and
they depend on a further half a gigabyte of mathematical libraries that are
present only on machines where a developer toolkit has been installed. Close to
a gigabyte, for one manufacturer's hardware.

The alternative had been overlooked because nobody had packaged it. The same
inference library can be built against a vendor-neutral graphics interface whose
runtime component is installed alongside every graphics driver. Querying that
interface on the development laptop, a machine with none of the manufacturer's
software installed, and which had been recorded in these notes as having an
unusable card, returned three devices, the integrated processor graphics, the
discrete card through an open-source driver, and a software renderer. Two of the
three were immediately usable. The proprietary route would have reached neither.

The download for that approach is about ninety megabytes and requires nothing
else, because the driver component is already present on any machine that can
display graphics at all. The trade is genuine: the vendor-neutral path is slower
on that vendor's hardware than the proprietary one. It is far faster than the
processor on every machine, which is the comparison that matters when the
alternative is no acceleration at all.

The pattern is worth recording because this is the third time it has appeared.
The language model is not included in the installer because of a file size limit
imposed by the distribution platform. Typesetting worked everywhere only once
the typesetting engine was included rather than assumed. And graphics hardware
was unusable because of which build was shipped. In each case the question
looked like ``is this machine capable'' and the answer was ``yes, but can it be
delivered''. Capability was never the constraint; packaging was.

Two errors were made during the work and are worth recording as much as the
result. The first: the software renderer must be excluded. It is the processor
presenting itself as a graphics device, so directing work to it would route
through a translation layer to reach the hardware already performing the
calculation, while reporting success at every step. The second: the memory each
device reports is not its own. Both real devices claimed twelve gigabytes, which
is the system's memory rather than the card's four, and the software renderer
claimed the most of all. An early version ranked devices by reported memory,
which would have selected precisely the device that must never be selected.
Both were found by running the detection and reading the output, not by
reasoning about it.

A third error was caught by the supervisor rather than by any test. The module
that answers hardware questions is deliberately constructed so that every fact
is supplied to it rather than discovered by it, which is what allows a
fabricated machine to be reasoned about on hardware that resembles it in no way.
A convenience function was added that consulted the real graphics interface from
inside that module. It was noticed immediately, from the outside, by someone
asking whether that module was not supposed to be self-contained. It was, and
the addition had quietly reintroduced exactly the class of assumption the module
had been written to eliminate.

\subsection{8 August: a check that tested for silence}

Three failures in one day turned out to share a shape, and the shape is worth
recording because it is not obvious until the third one.

The graphics engine described in the previous entry was finally built and
published. It failed twice before it succeeded, for unrelated reasons, and then
failed a third time against its own verification step. That step asserted that
the compiled library reported support for offloading work to the graphics
processor. It reported that it did not. The library was correct: ninety
megabytes of compiled shader code, exactly as intended. The function being
asked, however, reports whether a graphics backend has \emph{registered at
runtime}, not whether one was \emph{compiled in}, and the machine performing the
build has the development libraries but no graphics driver. The assertion was
testing the build machine rather than the artefact.

The same day, a release built correctly, uploaded every file, reported every
stage successful, and remained unpublished. An unpublished release is invisible;
requests for its files return "not found"; and the update mechanism treats "not
found" as "you are already up to date". Testers were therefore told they were
current while the release sat unreachable. The instruction to publish had been
given. It had simply not been obeyed, and nothing checked.

The third instance was the most instructive because it is still in the
application. Metadata extraction falls back to the language model when the
faster method fails, and the condition guarding that fallback tests whether the
extracted title is \emph{empty}. The faster method almost never returns an empty
title. It returns a wrong one. The better path is therefore unreachable, not
because it is broken, but because the broken path always reports success.

The common form: a guard that checks for absence, protecting against a failure
mode that is confident wrongness. Absence is easy to detect and rare. Confident
wrongness is common and looks identical to success. Where a value can be
legitimately empty, zero, false, an empty string, a check for emptiness
cannot distinguish "no answer" from "a bad answer", and the second is the one
that reaches the user.

\subsection{9 August: refusing to build on a layer I could not vouch for}

The next feature was citation support: allowing an author to cite a paper from
their own collection, which the application can do and comparable tools cannot,
because it has already read the paper. The entry it would generate is built from
metadata extracted at ingestion, and that extraction was written by a
collaborator. Rather than build on it, I asked to see it work.

The instrument was a short script that follows one document through every stage
and preserves the output of each, so the intermediate results could be examined
rather than described. Three published papers of established standing were used,
so that correct answers were known in advance.

Not one of the three produced a usable author list. On the best known of them
the extracted title was the copyright notice printed above it; on the others the
title was correct but had authors' names attached to the end of it. In every
case the author list contained institutions, the employers named beneath the
authors, and in two cases the paper's own title appeared as an author. The
publication year, taken as the first four-digit number appearing near the start
of the document, was wrong for the most significant of the three.

The cause is that the extractor receives the document as unstructured text. A
document of this kind is not text; it is positioned glyphs, each carrying a size
and an orientation, and the title of a paper is reliably the largest horizontal
text near the top of the first page. That information exists in the file and is
discarded before extraction begins. Restoring it, with two exceptions for a
rotated identifier stamp printed down the margin of one archive's documents,
identified all three titles exactly.

Two things are worth separating here. The first is that the fault was not one of
implementation quality but of information: a correct implementation of the wrong
question. The second is that I had considered rewriting the component in a
systems language for reliability, and measurement showed that the entire stage
accounts for under one percent of processing time, the remainder is the
embedding model, which is already compiled native code. A rewrite would have
reproduced the same wrong answer faster.

This metadata labels every node of the literature map. The map has therefore
been displaying copyright notices where paper titles were intended, in released
software, for some weeks. It was not reported by anyone.

\subsection{12 August: a measurement that contradicted its own test set}

The repair followed from the diagnosis. Extraction was rebuilt to read the
positioned glyphs rather than the flattened text, separated into two components:
one concerned only with page geometry, which groups glyphs into lines and divides
each line at its column gaps, and one concerned with what a title or an author
is. The separation was not tidiness. Three of the four most damaging faults were
geometric and had nothing to do with scholarly documents: lines grouped by the
top of their bounding box rather than by the baseline they share, which split a
title set in small capitals into single letters; a change of font within a word,
which inserted a space into it; and an accent stored as a separate glyph before
its letter, which left a surname unsearchable. None of these are visible while
reading a procedure that is simultaneously deciding whether something is a name.

Removing institutions from author lists was treated as a question of structure
rather than of vocabulary. A list of employers is never complete and dates
immediately. Three rules were used instead: a structural word such as
\emph{university} or \emph{institute} disqualifies an entire line rather than a
single entry, since a line of affiliations may name only some of them in terms a
list would recognise; a short capitalised abbreviation standing beside ordinary
words denotes an organisation; and a phrase printed more than once within the
author region is an address rather than a person, because a shared address is
printed once per author while a name is printed once per document. The third rule
requires no vocabulary of any kind, and is the only one capable of distinguishing
a country from a surname, the two being indistinguishable by spelling.

The result on the documents used during development was complete: every title,
year and author list exact. That number was then discovered to be worthless. A
second measurement was constructed in which neither the documents nor the
expected answers were chosen by the author: papers were sampled across fifteen
subject classifications through a public archive interface, and the archive's own
catalogue records were taken as the expected values. Accuracy on that sample was
sixty-two percent, against one hundred on the development set.

The entire difference lay in typesetting conventions that had not been examined.
One physics template spaces a centred author line widely enough that the spacing
is indistinguishable from column separation, so a line of two authors was read as
five fragments and none of them as a name. Another runs from the affiliations
directly into the abstract without printing a heading, so the search for authors
continued into the body of the paper and collected section headings. A margin
threshold, introduced to discard an archive's identifier stamp, was discarding
the first line of any centred title wide enough to begin left of it. Correcting
these raised the sample to ninety-three percent of titles and eighty-six percent
of author lists exactly correct.

The lesson is narrower than it first appears. The development set had not been
chosen dishonestly; it had been chosen from a single field, and a template is a
convention of a field. A measurement taken over material one has selected
oneself describes the selection, not the method. The remaining failures,
mathematics within a title, publisher cover pages preceding the document, and
submissions too recent to carry an identifier: are a long tail of the same
kind, and the established solution to a long tail of layout conventions is a
classifier trained over layout features rather than a further accumulation of
rules. That is what the mature tool in this area does, and it is the point at
which this approach should be expected to stop improving.

\subsection{13 August: an interface reporting what it had inferred}

The collection screen is one member's work throughout, and it had been extended
with panels summarising the state of the application: which model was in use,
how many analyses had run, what was being written. Reviewing the layers beneath
those panels found that two of the four figures had never displayed a value in
any release, they were fixed placeholders, and that a third was inferred
rather than read.

The inference is the instructive one. The panel naming the model in use
searched the list of installed models for the first whose status was
\emph{ready}. No such status exists; the application reports \emph{present}.
The search therefore never matched, and the panel silently displayed whichever
model happened to be first in the list. Beneath that lay a larger error: there
is no single model in use. Work is routed by task, and which model serves a
task depends on every other model installed: what it has been assigned, how
much memory is free at that moment, and precedence between them. The backend
computes that and records the reason in a comment written weeks earlier:
duplicating the rule in the interface would let the two drift. On the
development machine the larger model that actually performs analysis did not
appear in the list being searched at all. The panel now reads the routing table
the backend already produces.

A related fault ran through the same screen, a single dash stood for three
different facts, not yet loaded, nothing to report, and the request failed.
An outage was therefore indistinguishable from an empty library. This is the
same shape as the extraction guard described above, and the third instance
recorded this month of one value being asked to carry several meanings.

Two further findings concerned the layout rather than the data. Both of the
pane-based screens sized themselves by subtracting a constant from the height
of the window, with a comment estimating the height of the page heading. That
is a measurement replaced by a guess: when the headings were removed the
constant was wrong, silently, and the panes stopped reaching the bottom of the
window. Replacing it with a layout that fills the space it is given failed on
the first attempt, because the property used to do so resolves against the
nearest enclosing flexible container and an unclassed wrapper sat between the
two. Both screens then collapsed to the height of their own contents. The
correction was to pass the height explicitly at every level.

The lesson recorded is not about the technique. A mechanism was replaced
without first establishing what it sat inside; the original was fragile, but it
worked, and the replacement was correct in principle and wrong in that
structure. The same mistake recurred twice more in the same session,
filling a column whose neighbour then grew taller, and equalising panels whose
contents were not equal: before the layout was expressed as flow rather than
as position.

\subsection{14 August: a scale, and a sidebar that leaves a rail}

The author of the collection screen judged that it presented too much at once,
and reduced it: six summary figures that had opened the page were moved to
caption the section they describe, two were removed outright, and the document
list was divided into pages. The removals are the interesting part. One
reported the number of bytes the stored papers occupy: accurate, cheap to
compute, and of no use to any reader, since nobody is short of the space in
question and the figure changes nothing about what they do next. The other
reported how many fragments the text had been divided into for indexing, which
describes the implementation rather than the library. A figure earns its place
by answering a question the reader would actually ask, and neither did.

Two faults found during that review belonged to the shell rather than to the
screen.

The first was that the interface had twelve different small type sizes,
0.68, 0.70, 0.72, 0.75, 0.76, 0.78, 0.80, 0.82, 0.84, 0.85, 0.90 and 0.95 of
the root size: distributed across a hundred and sixty-one declarations in
two stylesheets. Differences of two hundredths are not decisions; they are what
accumulates when each component chooses independently, and the effect is an
interface that reads as careless without any single value being identifiably
wrong. These were consolidated to six named steps. The root size was also made
to vary with the width of the window rather than being fixed, since the
application is read at arm's length on displays between thirteen and
twenty-seven inches. A first attempt allowed too wide a range and produced text
that read as magnified rather than as legible; the range was narrowed.

The second concerned the navigation panel, which collapsed to nothing. That
sounds like the greatest possible gain of space and is the reverse, the two
controls still required when it is closed, the mark that reopens it, and the
button explaining the current screen, then have to be drawn on top of the
page, obscuring whatever is beneath them. Collapsing the panel to get it out of
the way therefore put two things in the way. It now collapses to a narrow rail
that carries both, and because the page's margin is derived from the same
declared width, the content moves aside for the rail rather than passing under
it.

One defect from that change is worth recording for its shape. Moving the
explanatory button onto the rail made it vanish. The panel is drawn above it in
the stacking order, so the button was rendered behind the rail: present,
correctly positioned, and invisible. Nothing in the change suggested a stacking
question, and nothing in the source indicated one; it was found by looking at
the screen and observing that the control was not there. It is a reminder that
a layout change is verified by looking, not by reasoning about what should
follow.

\subsection{14 August: replacing the interface}

The interface was replaced. The original had reached a point where the same
class of correction kept recurring: a dozen font sizes chosen independently by
different components, a navigation panel that collapsed to nothing and
therefore had to draw its remaining controls on top of the page, and panes
sized by subtracting an estimated constant from the height of the window. The
decision taken was to replace it rather than continue correcting it.

All three members worked on the replacement. The design was built as a
parallel tree while the original continued to ship, so that neither the
released application nor the work in progress blocked the other.

The arrangement has one failure mode, and it governed how the work was done: two
trees drift, and the one nobody is looking at rots. Three measures were taken
against it. Both trees were run through the same continuous integration matrix,
so neither could reach a release having been checked differently. The test that
verifies the interface calls only routes the backend defines was found to name
the old directory explicitly, which would have meant that on the day the
replacement shipped it went on passing while examining a directory nobody
built; it now takes the directory as a parameter. And a third check was added
that calls every route against a running server, because parsing both sides
proves the paths agree while only calling them proves anything executes.

Every correction made during the fortnight both trees existed was made twice.
That is the true cost of replacing a running interface, and it is not the
rewrite but the period of double maintenance around it. Copying between the
trees was not available: the replacement uses the token names the original used
for sizes to name its colours instead, so a mechanical transfer would have
produced a stylesheet in which a colour answered to a size.

One defect from that period is worth recording for its shape. Moving the
button that explains the current screen onto the collapsed rail made it
disappear. It was present, correctly positioned and invisible, because the
panel is drawn above it in the stacking order. Nothing in the change suggested
a stacking question and nothing in the source indicated one; it was found by
looking at the screen and observing that the control was not there.

\subsection{14 August: citations, and two failures that printed instead of raising}

The paper editor could compile a document but could not cite one, which for a
research tool is the omission that matters most. The interaction chosen was
deliberately not a dialog: an author is already typing when they decide to cite,
so typing the word \texttt{cite} opens the library beneath the cursor, further
typing filters it, and pressing Enter replaces the word with a citation command.
Choosing nothing leaves the word where it was, and an ordinary word compiles, so
nothing is committed by opening the list.

One planned subsystem was removed before it was built. The intention had been to
track citation numbers so that the numbering stayed correct as citations moved.
BibTeX already renumbers every citation on each recompile and orders the
reference list according to the chosen style, so the entire mechanism was
available at the cost of emitting keys instead of numbers. It is worth recording
as a decision because the alternative would have been a second, worse copy of
something the toolchain performs correctly.

Two defects were found underneath, and they share a shape with the extraction
defect recorded on 12 August. Neither raised an error.

The first was in storage. Author lists were flattened into the vector store by
joining them with commas, because the store holds only strings. In a
bibliography the comma is not a separator between people, it separates a surname
from a given name, so a list of eight authors read back as sixteen fragments of
names. Nothing failed. A bibliography does not validate its input; it typesets
it. The list is now encoded as structured text, and the older form is still read
so that a library already built does not have to be rebuilt.

The second was that a document containing citations but no bibliography command
produced neither the references nor an error. The command that prints a
reference list is also the command that instructs the compiler to run the
bibliography processor at all, so its absence removed both the output and the
process that would have produced it. Each citation was typeset as a question
mark. The compile step now supplies the command when a document cites something
and the project has a bibliography file to point at, both conditions being
necessary: naming a file that does not exist converts a working compile into a
failing one.

A third defect is recorded for the method rather than the cause. The citation
list did not appear. Every automated test passed, including tests that typed
into an editor and asserted that the list opened, because those tests run in a
simulated document object model that computes no layout. Driving the actual
application and reading the element's position back showed that the list was
present, populated with every paper in the library, and positioned 1382 pixels
down a window 1000 pixels tall, inside a region that hides rather than scrolls
whatever exceeds it. A second cause was found the same way: the compiled
document displayed beside the editor was drawn above the list, concealing the
column that held the citation keys.

The lesson is narrow and worth stating exactly. A test environment that omits
layout cannot fail a layout defect. A feature correct in every respect except
where it is drawn is, to the person using it, a feature that was never built.
What this means for the interface suite is set out in the testing document.

\subsection{14 August: repairing what extraction gets wrong}

Extraction identifies titles correctly for about 93 percent of papers and author
lists for about 86 percent. Roughly one paper in seven is therefore stored under
something wrong, and that record names the paper on the literature map, in every
search result, and in every reference generated from it. A citation feature built
on top of it inherits the error.

The question was where the correction should be made, and the two candidates
were not equivalent. Editing a document's bibliography repairs that document;
the next one is generated from the same wrong record and repeats the mistake.
Editing the library record repairs the source, and every citation, map node and
search result made afterwards is correct. The library was chosen for that
reason, and the editor was deliberately left without a second facility of its
own, since a project's bibliography is already a file in its tree and can be
opened like any other.

One detail is worth recording because it was nearly got wrong twice in the same
day. The editing form separates author names on commas, but accepts semicolons
instead when any are present, because a name written surname first contains a
comma. Splitting on it is precisely the defect the morning had been spent
removing from the storage layer.

\subsection{15 August: what other people's machines found}

Version 2.3.0 reached the stable channel and, for the first time, testers on
macOS and Windows ran the replaced interface and the writing feature on
hardware nobody here owns. The round produced four defects in an afternoon,
after a fortnight in which every automated suite had been green.

None of the four was a regression. All had been present in earlier builds and
none had been observed, which is the useful part of the result: the suites were
not wrong, they were looking at the wrong layer.

Two are worth recording for their shape.

On Windows the reader could not display a paper, and the message named a file
rather than a document, so the tester reasonably deleted the papers and ingested
them again. Nothing changed, because no paper was involved. Python assembles its
table of media types partly from the Windows registry, which has no entry for
the extension the viewer's worker module uses; the server therefore described
that module as plain text, and a browser refuses to execute plain text as a
program. The file arrived intact and was declined on the doorstep. On Linux and
macOS the built-in table answers correctly, so nothing in development or in the
test matrix could have seen it: the defect is a browser's reaction to a header,
and no test in the suite reads headers.

The second was raised independently by testers on two different operating
systems within an hour of each other: selecting a point on the map takes more
precision than it should. When two people on different machines describe the
same difficulty without prompting, the cause is the design rather than the
hardware. It was geometry. A paper is drawn as a point a few units across and
the click target was the drawing itself, which is about six pixels once the map
is scaled to show everything. The drawing was left alone, since larger discs
were removed months ago for turning the map into a bubble chart, and an
invisible target now carries the selection at a constant size on screen however
far the view is zoomed. The label was also unselectable, which mattered more:
it is thirteen times the area of the point, and a reader who has found the
paper they want has almost always found it by reading the title.

A third finding was not a defect at all. A tester said the reading surface was
tiring in a dim room, which is a physical report rather than a preference: the
page was at 94\% luminance and filling the screen. Lowering it turned out to
have a trap in it. An intermediate version lowered the page and the surface
beneath it together and was rejected on sight as looking dusty, because
luminance had fallen while colour intensity rose, and dark with strong colour
reads as dirt rather than as paper. The final version lowers only the page, and
puts the character of the stock into its warmth instead of into its darkness.

\section{Division of Work}
\markboth{Division of Work}{}

Work was divided by feature rather than by layer, so that each member owned a
capability end to end: its data model, its interface, and its failure modes
rather than a horizontal slice of everybody's. The consequence is that no
part of the product can be described without naming the person who built it.

\begin{center}
\small
\begin{tabularx}{\linewidth}{@{}llX@{}}
\toprule
\textbf{Capability} & \textbf{Principal author} & \textbf{Scope} \\
\midrule
Retrieval and knowledge base & Aditya & document ingestion, chunking, embedding
index, semantic search \\
Analysis & Aditya & summarisation, claim extraction, theme clustering, and the
precomputation that moved them off the request path \\
Gap finder & Aditya & gap detection over the analysed corpus, and the
suggestion engine above it \\
LitGraph & Aditya & the canvas consolidating search, analysis and gaps; the
reader, passage marking and node interaction \\
\midrule
Scribe & Rithesh & LaTeX editor, the auto-healing compiler, and the natural
language to markup path \\
Desktop shell & Rithesh & the Tauri application, native hardware diagnosis,
process supervision, in-application updates \\
Delivery engineering & Rithesh & cross-platform build, release pipeline,
artefact validation, and the automated test suite \\
Model management & Rithesh & hardware capability model, task routing, the
model registry, and Hugging Face acquisition \\
\midrule
Library & Jitvan & the paper collection screen in full, the document list,
its expansion and deletion, the encryption controls, and the overview panels
reporting collection size, the models in use and the drafts in progress \\
Encryption & Jitvan & at-rest encryption and decryption of stored documents,
with the key handling and interface around it \\
Distribution site & Jitvan & the public landing and download page \\
Interface refinement & Jitvan & typography and identity, dark mode correction,
and the assistant panel \\
\bottomrule
\end{tabularx}
\end{center}

The boundaries were not absolute. The interface was worked on by all three, and
several features crossed hands: encryption was introduced by Jitvan and later
extended when the key derivation was hardened; the desktop shell was begun as
packaging work and grew into hardware diagnosis, which then fed the model
routing. Where a capability changed owner, the original author's design was
retained and the reasoning recorded in the decision log rather than replaced
silently.

Review was mutual and enforced by the branch rules described above: no change
reached the release branches without passing another member's review and the
automated gate, including changes by the member who wrote those rules.


\section{Present State}

\begin{center}
\small
\begin{tabular}{@{}lr@{}}
\toprule
\textbf{Measure} & \textbf{Value} \\
\midrule
Commits & 347 \\
Pull requests merged & 86 \\
Application code (Python) & 14\,438 lines, 83 modules \\
Test code & 8\,375 lines, 48 modules, 1\,085 tests \\
Interface tests & 222 tests, including every screen mounted \\
Interface code & 13\,930 lines, 23 components \\
Desktop shell (Rust) & 777 lines \\
Build and release tooling & 4\,484 lines, 22 scripts, 7 workflows \\
REST endpoints & 63 \\
Architecture decisions recorded & 78 \\
Releases published & 13, most recently v2.3.0 \\
Platforms built and validated automatically & 3 \\
\bottomrule
\end{tabular}
\end{center}

Figures are counted on the working branch at 15 August and include the citation
work described above, which is complete and tested but not yet released. The
totals are of tracked files only: generated output, dependencies and the local
evaluation scripts are excluded, so the counts describe what the team wrote
rather than what the repository contains.

\subsection{Verification status by platform}

\begin{center}
\small
\begin{tabularx}{\linewidth}{@{}lX@{}}
\toprule
\textbf{Platform} & \textbf{Status} \\
\midrule
Linux & fully verified. The development platform: every suite runs here, and the replaced interface, the paper editor, citations and the map are exercised by hand on physical hardware, including against 2.3.0 \\
Windows & verified. 2.3.0 installed and exercised by a tester: writing, citing and compiling confirmed working. The reader defect of 15 August was found here and is fixed \\
macOS & verified. 2.3.0 installed and exercised by a tester: writing, citing and compiling confirmed working, including the citation list. Unsigned applications are still obstructed on first launch, which is a signing matter rather than a defect \\
In-application update & working; testers update in place from the beta channel \\
\bottomrule
\end{tabularx}
\end{center}

\section{Planned Work}

The following are planned rather than aspirational. Each follows from a
limitation identified during development, and several are already partly
prepared for in the existing structure of the system.

\subsection{Language models trained for this application}

The clearest quality limitation is the capability of a general-purpose 0.5
billion parameter model on structured output. Asked to plot a function it
produced a reference to an image file that does not exist, where a 1.5 billion
parameter model produced a correct plot. The intended response is models suited
to the task rather than simply larger ones, in three stages of increasing
ambition.

\paragraph{Stage one: fine-tuned task models.} Training pairs of natural
language instruction and compilable LaTeX are already collected passively during
normal use, and the routing table in the inference client already contains
entries for a dedicated writing model, so adopting one is a matter of supplying
weights. A second model fine-tuned for the structured summarisation and claim
extraction formats the application parses would address the cases where a
general model returns prose where JSON was requested.

\paragraph{Stage two: adapter distribution.} Rather than shipping whole models,
distribute small adapters applied over the bundled base model, so a capability
upgrade costs tens of megabytes instead of gigabytes. Selection would use the
hardware profile the application already computes at startup, so a machine
receives only variants it can run.

\paragraph{Stage three: models trained from scratch.} The longer term intention
is to train small language models specifically for academic literature work,
rather than adapting general-purpose ones. A model trained from the outset on
research text, structured extraction and typesetting markup can be far smaller
than a general model of equivalent usefulness for these tasks, which matters
directly for a product constrained to consumer hardware.

\paragraph{Federated learning.} A possible direction, subject to careful design,
is improving these models from real usage without collecting anyone's documents.
Training would occur on the user's own machine and only model updates, never
text, would leave the device, and only with explicit consent. This is recorded
as a possibility rather than a commitment: it must not be allowed to weaken the
privacy guarantee that is the reason the product exists, and that constraint
takes precedence over the capability.

\subsection{An interface design philosophy, and the refactor that follows it}

The interface was built feature by feature under time pressure. It is coherent
but was not designed as a system: styling has accumulated into a single large
stylesheet with two competing styling mechanisms in use.

The intention is deliberately ordered. A design philosophy will first be agreed
by the team through discussion, and the interface will then be refactored to it.
Refactoring before that agreement would simply produce a second set of
accumulated decisions. Points to be settled include:

\begin{itemize}[leftmargin=*]
  \item \textbf{Design tokens.} Colour, spacing and typography expressed as
        variables, so the interface compiles from a defined system rather than
        from decisions made one screen at a time.
  \item \textbf{One environment rather than four.} Browsing, searching,
        interpreting and writing are currently separate screens. The intended
        direction is a single workspace whose density and visual quiet change
        with the task.
  \item \textbf{A single cross-platform identity} rather than per-platform
        conventions, so the application looks deliberate on every operating
        system.
  \item \textbf{Latency as a design material.} Local inference takes seconds and
        that cannot be removed. The interface should be designed around honest
        waiting rather than concealing it.
  \item \textbf{Reading comfort.} Sessions are long and the content is dense, so
        typography, measure and contrast are functional requirements rather than
        decoration.
\end{itemize}

\subsection{A reliable paper editor}

The paper writer is the most used feature and has received the most correction,
but its editing surface remains a plain text area over LaTeX source. The planned
work is a genuinely capable editor:

\begin{itemize}[leftmargin=*]
  \item syntax awareness, bracket matching, and compiler diagnostics shown
        against the line that caused them;
  \item a visual editing mode in which document structure can be edited without
        the author reading markup, with the compiled PDF remaining the single
        source of truth for appearance;
  \item position synchronisation between the source and the compiled document,
        so that selecting a place in one moves to it in the other;
  \item continued strengthening of the natural language to LaTeX path, which
        currently rewrites a selected passage in place and is the feature the
        planned writing model would most improve.
\end{itemize}

\subsection{Further planned refactors}

\begin{itemize}[leftmargin=*]
  \item \textbf{Incremental keyword index}, removing the per-query rebuild that
        currently limits corpus size.
  \item \textbf{Stylesheet consolidation}, resolving the two styling mechanisms
        into one, carried out as part of the interface refactor above rather
        than separately.
  \item \textbf{Code signing and notarisation} for macOS and Windows. This is
        required before a production release: both operating systems currently
        obstruct the first launch of an unsigned application, which every user
        encounters and which reads as the software being broken.
  \item \textbf{Bundled model integrity checks}, so that a corrupted or partial
        payload is detected at startup rather than at first use.
\end{itemize}

\section{Closing Note}

The technically demanding parts of this project, namely local inference with
hardware-aware model selection, hybrid retrieval, an auto-healing typesetting
pipeline and authenticated encryption, were completed and work. What determined
whether anyone else could use the software was packaging, dependency
completeness, and verification of the delivered artefact.

The distinction between software that works and software that has been shown to
work on a machine other than its author's is the lesson this project taught most
clearly, and it was taught by a release that failed rather than by one that
succeeded.

\end{document}
